\section{The complete observable table constrains an instrument analysis}
\label{exec17:s3-sensitivity}

The instrument example above leaves a practical question unanswered. When the
bank doubts a premise, which conclusions survive, and what additional
information would narrow the answer? A Wald ratio alone cannot answer this.
It discards restrictions in the joint receipt--outcome distribution, while
an assumption change alters the compatible population. We develop the full
binary problem before making either change. Every probability in the following
construction is authored; the sensitivity budgets express hypothetical
scientific judgments rather than facts recovered from bank records.
\citep{exec17_s3_sensitivity}

Let $Z$ be a randomized encouragement, $D(z)$ receipt under encouragement $z$,
and $Y(d)$ success under receipt $d$. For now retain consistency, exclusion
and an encouragement independent of all potential responses, with positive
probabilities for both assignments. Exclusion makes the outcome under the
realized encouragement equal to $Y(D(Z))$. A response type is
$t=(d_0,d_1,y_0,y_1)\in\{0,1\}^4$ with mass $p_t$. The observable
conditional cell probabilities satisfy
\begin{align}
 q_{zdy}
 &=P(D=d,Y=y\mid Z=z) \notag\\
 &=\sum_{t\in\{0,1\}^4}
      p_t\,{\bf1}\{d_z=d,\ y_d=y\}, \qquad
 p_t\geq0,\quad \sum_t p_t=1 .
 \label{exec17:s3-response-map}
\end{align}
The same sixteen masses must generate both assignment groups. Fitting two
unrelated receipt--outcome tables would lose the common-population
restriction supplied by randomization. The assignment probability does not
multiply these conditional cells; it enters the unconditional law and the
sampling precision.

The population receipt effect is the linear objective
$\sum_t p_t(y_1-y_0)$. Its sharp bounds under these assumptions are the
minimum and maximum of this objective subject to
(\ref{exec17:s3-response-map}). This finite optimization represents the
stated information class exactly: every admissible potential-response law
has these masses, and every feasible mass vector constructs an admissible
law by drawing the type independently of $Z$. A constraint on defier mass
is also linear. In contrast, the effect among compliers, the types with
$d_0=0,d_1=1$, is
\[
 \theta_C=
 \frac{\sum_{t:d_0=0,d_1=1}p_t(y_1-y_0)}
      {\sum_{t:d_0=0,d_1=1}p_t}.
\]
This is a linear-fractional target with a scientifically meaningful
denominator. A zero denominator means there is no such stratum; a
numerical convention cannot define its mean effect.

Use the following conditional table, which retains all receipt and outcome
cells of the earlier example:
\[
\begin{array}{c|rrrr}
 &D=0,Y=0&D=0,Y=1&D=1,Y=0&D=1,Y=1\\ \hline
 Z=0&13/25&9/50&9/50&3/25\\
 Z=1&6/25&3/50&19/50&8/25
\end{array}
\]
Each row sums to one. Receipt means are $3/10$ and $7/10$; outcome
means are $3/10$ and $19/50$. Thus the receipt first stage is
$d=2/5$ and the outcome intention-to-treat contrast is
$\tau_Z=2/25$. Their ratio is $1/5$. These two contrasts are summaries
of eight cells, rather than a substitute for them.

\section{Sharp population bounds with monotone receipt}
\label{exec17:s3-monotone-polytope}

Monotonicity, $D(1)\geq D(0)$ for every customer, removes the four
defier response types. Let $N,C,A$ denote never-receivers, compliers and
always-receivers. Their masses are determined by the receipt means:
\[
 \pi_N=3/10,\qquad \pi_C=2/5,\qquad \pi_A=3/10.
\]
Under $Z=1,D=0$ the observed group consists only of never-receivers;
its success mass is $3/50$. Under $Z=0,D=1$ it consists only of
always-receivers; its success mass is $3/25$. Subtracting those
contributions from the other two success cells gives
\[
\begin{array}{c|cc}
 &E\{Y(0)\mid\text{stratum}\}&E\{Y(1)\mid\text{stratum}\}\\ \hline
 N&1/5&n\\
 C&3/10&1/2\\
 A&b&2/5
\end{array}
\qquad 0\leq n,b\leq1 .
\]
Here $n$ and $b$ are conditional means. Below, when principal-stratum
masses vary, we will use different symbols for unconditional success
masses. The distinction prevents multiplication by a stratum weight twice.

The complier effect is $1/2-3/10=1/5$, as the LATE argument predicts
under independence, exclusion, relevance and
monotonicity.\citep{angrist1996}
The population receipt effect instead is
\begin{align}
 \theta
 &=\tfrac{3}{10}(n-\tfrac15)
   +\tfrac25(\tfrac12-\tfrac3{10})
   +\tfrac3{10}(\tfrac25-b) \notag\\
 &=\tfrac7{50}+\tfrac3{10}(n-b),
 \qquad
 -\tfrac4{25}\leq\theta\leq\tfrac{11}{25}.
 \label{exec17:s3-monotone-ate}
\end{align}
The table supplies neither the outcome of giving receipt to never-receivers
nor the outcome of withholding it from always-receivers.
These are missing counterfactual responses, even with infinite observational
precision and perfect randomization of encouragement.

It remains to prove attainability rather than merely bound a formula.
Within a stratum of mass $m>0$, let $s_0,s_1$ be its success masses
under receipt zero and one, each in $[0,m]$. One admissible coupling is
\[
 p_s(y_0,y_1)=
 \frac{a_0(y_0)a_1(y_1)}{m},\qquad
 a_j(1)=s_j,\quad a_j(0)=m-s_j .
 \label{exec17:s3-product-coupling}
\]
Its four nonnegative cells sum to $m$ and have the required two margins.
Zero-mass strata contribute four zeros. Conditional independence of the two
responses is used here only to construct a witness; it is not imposed on
every model in the bounded class. Other couplings are allowed.

Choose $(n,b)=(0,1)$ for the lower endpoint and $(1,0)$ for the upper
endpoint, and apply (\ref{exec17:s3-product-coupling}) in each of $N,C,A$.
Substituting the resulting twelve possibly nonzero type masses into
(\ref{exec17:s3-response-map}) reproduces every cell of the displayed
table. Convex mixtures of the two type laws reproduce that same table
and attain every intermediate population effect. Equation
(\ref{exec17:s3-monotone-ate}) is consequently a sharp range for this
class. At the stipulated USD 100 per success and USD 12 per receipt,
universal receipt has incremental net range $[-28,32]$ USD per eligible
customer. A positive complier effect therefore does not settle this
population-wide decision.\citep{exec17_s3_sensitivity}

\section{A bounded departure from monotonicity}
\label{exec17:s3-defiers}

Suppose some customers may refuse the service because the encouragement
was sent. Retain randomization and exclusion, but stipulate
\[
 f=P\{D(0)=1,D(1)=0\}\in[0,1/10].
\]
The observed first stage identifies $\pi_C-f=2/5$, rather than
$\pi_C$ itself. Consequently
\[
 \pi_N=\pi_A=3/10-f,\qquad \pi_C=2/5+f.
\]
The $1/10$ budget is an assumed substantive restriction. Neither the
observed first stage nor an absence of obvious complaints verifies it.
The admitted complier population now varies across compatible models;
it should not be described as a fixed set of known customers.

Let $h_j=P(\text{defier},Y(j)=1)$ be a defier success mass. Subtracting
principal-stratum contributions from all four observed success cells gives
\[
\begin{array}{c|cc}
 &P(\text{stratum},Y(0)=1)&P(\text{stratum},Y(1)=1)\\ \hline
 N&3/50-h_0&u_N\\
 C&3/25+h_0&1/5+h_1\\
 \text{defier}&h_0&h_1\\
 A&u_A&3/25-h_1
\end{array}
\]
where the two unobserved masses satisfy
$0\leq u_N,u_A\leq3/10-f$. Nonnegativity and the stratum sizes imply,
for the declared range of $f$,
\[
 0\leq h_0\leq\min(f,3/50),\qquad 0\leq h_1\leq f .
 \label{exec17:s3-defier-feasible}
\]
There are no further constraints on these margins for this example.
For instance, the never-receiver $Y(0)$ mass is at most its stratum
mass since $3/50\leq3/10-f$, and the always-receiver $Y(1)$ mass
is nonnegative since $h_1\leq1/10<3/25$.
The complier success masses are below $2/5+f$ throughout.
Applying (\ref{exec17:s3-product-coupling}) separately to all four
strata then proves sufficiency. It also reproduces the failure cells,
since the receipt totals and success cells agree.

The complier target becomes
\[
 \theta_C(f,h_0,h_1)
 =\frac{2/25+h_1-h_0}{2/5+f}.
 \label{exec17:s3-defier-ratio}
\]
For fixed $f$ it is increasing in $h_1$ and decreasing in $h_0$.
Its smallest feasible numerator uses $h_1=0$ and
$h_0=\min(f,3/50)$. On $0\leq f\leq3/50$ the lower ratio is
$(2/25-f)/(2/5+f)$, with strictly negative derivative.
On $3/50\leq f\leq1/10$ it is
$(1/50)/(2/5+f)$, again decreasing.
The maximum uses $h_0=0,h_1=f$ and
$(2/25+f)/(2/5+f)$, whose derivative is positive.
Both extrema occur at $f=1/10$, yielding
\[
 \boxed{\ \theta_C\in[1/25,9/25]\ },\qquad
 100\theta_C-12\in[-8,24]\ {\rm USD}.
 \label{exec17:s3-defier-sharp}
\]
The lower witness has $(h_0,h_1)=(3/50,0)$; the upper has
$(h_0,h_1)=(0,1/10)$. Their complier mass is the same, $1/2$.
Product couplings give complete nonnegative response-type laws at both
endpoints; mixtures preserve the table and this denominator, attaining
every intermediate complier effect.

Keeping only $d$ and $\tau_Z$ would give a weaker answer. The contrast
identity under exclusion is
\[
 \tau_Z=(2/5+f)\theta_C-f\theta_{\rm defier}.
\]
Allowing any defier mean effect in $[-1,1]$ gives the outer range
$[-1/25,9/25]$ for the same budget. Its negative lower endpoint
would require $f=1/10$ and $h_0=1/10$, which violates the
observed-cell restriction $h_0\leq3/50$.
Thus fitting a Wald decomposition is insufficient to certify full-law
compatibility. The complete table rules out that negative effect, while
still leaving the USD 12 receipt decision unresolved.

The population effect has a different response to this relaxation:
\[
 \theta=\tfrac7{50}+h_1-h_0+u_N-u_A.
\]
For any admitted $f$, its lower bound is
$-4/25+f-\min(f,3/50)\geq-4/25$ and its upper bound is
$11/25-f+h_1-h_0\leq11/25$.
The monotone endpoint laws remain admitted at $f=0$, so the sharp
population range stays $[-4/25,11/25]$. The larger model class
widens uncertainty for the complier target without widening this
population range. Sensitivity is a property of a specified target and
information class, rather than a single discount applied to every
reported effect.\citep{exec17_s3_sensitivity}

\section{Exclusion sensitivity needs a different response model}
\label{exec17:s3-exclusion}

A bank message may affect confidence or repayment even when the customer
does not receive the proposed counseling. Exclusion then fails.
The notation $Y(d)$ hides that pathway; use $Y(z,d)$ instead.
In this section restore monotone receipt and retain randomized encouragement,
but allow direct encouragement effects. Define
\[
 \theta_{C0}=E\{Y(0,1)-Y(0,0)\mid C\},\qquad
 \Gamma=E\{Y(1,D(1))-Y(0,D(1))\}.
\]
The first target changes receipt while holding encouragement at zero,
among the receipt compliers. The second averages a direct encouragement
contrast while fixing each customer's receipt to its $Z=1$ value.
It averages over all principal strata. A budget on $\Gamma$ is therefore
not automatically the same as a bound on each individual's direct effect.

Add and subtract $Y(0,D(1))$ inside the randomized contrast:
\begin{align}
 \tau_Z
 &=E\{Y(1,D(1))-Y(0,D(0))\} \notag\\
 &=\Gamma+
   E\{Y(0,D(1))-Y(0,D(0))\} \notag\\
 &=\Gamma+d\theta_{C0}.
 \label{exec17:s3-exclusion-decomposition}
\end{align}
The last equality uses monotonicity to make the receipt-changing group
precisely the compliers. It does not require exclusion.
If the scientific judgment is $|\Gamma|\leq g$, the implication is
\[
 \theta_{C0}\in
 \left[\frac{\tau_Z-g}{d},\,\frac{\tau_Z+g}{d}\right]\cap[-1,1].
 \label{exec17:s3-exclusion-outer}
\]
This is an outer sensitivity envelope. The derivation uses a necessary
contrast equation; it does not prove that every point can satisfy all
eight observed cells in the enlarged $Y(z,d)$ model.
Solving that larger response-type feasibility problem would be required
before calling the envelope sharp.

For the authored table, the implications are
\[
\begin{array}{c|cc}
 g&\theta_{C0}\text{ outer envelope}
   &100\theta_{C0}-12\text{ outer envelope, USD}\\ \hline
 1/50&[3/20,1/4]&[3,13]\\
 1/25&[1/10,3/10]&[-2,18]
\end{array}
\]
Under the smaller budget every compatible receipt effect would exceed
the stipulated cost threshold. Under the larger budget this necessary
envelope cannot settle the decision. It also does not prove that both
signs of net value are attainable in the complete response model.
The budgets are economically consequential assumptions: doubling a
permitted direct contribution changes what this calculation can recommend.
Neither excellent predictive fit nor a small standard error on the
Wald ratio validates those budgets.

\section{A feasible encouragement decision and the next experiment}
\label{exec17:s3-sensitivity-decision}

There is an action whose value uses less causal structure. Suppose the
bank can randomize the encouragement, value every success at USD 100,
pay USD 12 for every receipt, and pay an additional USD 2 when
encouragement is sent. Its incremental expected net from encouragement
one versus zero is
\[
 100\{E(Y\mid Z=1)-E(Y\mid Z=0)\}
 -12\{E(D\mid Z=1)-E(D\mid Z=0)\}-2
 =\tfrac65\ {\rm USD}.
\]
Randomization and consistency for these well-defined assignment policies
identify the two mean contrasts. This calculation remains valid in
the presence of direct encouragement effects or defiers; it does not
divide by the receipt first stage. It concerns offering encouragement,
rather than assigning receipt to unknown compliers. Its positive value
in this exact law is compatible with the unresolved population receipt
range and the sensitive complier comparison.

The information needs now follow from the particular uncertainty.
Recording receipt and outcome jointly preserves the full-cell restrictions
that sharpened the defier calculation. Randomizing encouragement more
precisely estimates those cells but does not observe both receipt
responses for the same person. To narrow population receipt bounds, an
ethically feasible study would need information on responses currently
missing in the never- and always-receiver strata, or defensible additional
restrictions. A design that varies message content separately from service
availability could address the direct pathway, but only if its delivery
and exclusion claims are justified for the new target. Repeating the
same encouragement experiment cannot by itself certify its own exclusion
restriction.

These exact-law results also do not supply sampling confidence intervals.
With estimated cell probabilities, the equality-constrained feasible set
may be empty because of sampling noise even when the population law is
compatible. A proposed inferential extension could form a justified joint
confidence region for the entire conditional table, then optimize over
both that region and admissible response laws. If the region contains
the true table and the assumptions hold, the optimized set contains its
compatible effects. Establishing the region's coverage requires the actual
assignment, clustering and missingness law. Optimizing over eight unrelated
point estimates, or treating a solver's feasibility tolerance as statistical
uncertainty, does not establish that guarantee.

For a small prototype, the full response-type map offers a concrete oracle:
check all cells, distinguish an empty compatible class from a wide effect
set, retain the positive ratio denominator, and return feasible extremal
laws with their bounds. The authored reference evaluates thirty-six
piecewise-domain vertices and separately reconstructs every observable
cell by principal-stratum mixing. Analytic derivatives establish that
the ratio extrema occur among those vertices. This verifies the finite
construction and its attainable ranges; it is a proposed comparison
standard for a future implementation and a basis for choosing a useful
public or generated experiment.\citep{exec17_s3_sensitivity}
\label{exec17:s3-sensitivity-end}
